\section{Introduccción}
El presente trabajo práctico tiene como objetivo familiarizarse con los instrumentos que permiten generar y medir señales de tensión, en este proyecto se trabajó con el generador de señales y el osciloscopio. También se busca aplicar conocimientos teóricos de respuesta transitoria y analizar el efecto de variar los valores de la resistencia, inductancia y capacitancia de un circuito de segundo orden en su respuesta transitoria.

\subsection*{Circuito de primer orden (RL)}
Para el circuito RL, la respuesta temporal de la tensión en la resistencia frente a un escalón de entrada se define como:
\[ v_R(t) = V_S \cdot (1 - e^{-t/\tau}) \]

Donde $\tau$ es la constante de tiempo del circuito:
\[ \tau = \frac{L}{R_{eq}} \]

Donde la resistencia equivalente debe considerar las no idealidades de los componentes e instrumentos:
\[ R_{eq} = R_{gen} + R_f + R_L \]

\subsection*{Circuito de segundo orden (RLC)}
En el circuito RLC serie, el comportamiento transitorio está regido por la siguiente ecuación característica:
\[ s^2 + \frac{R_{eq}}{L}s + \frac{1}{LC} = 0 \]

Las raíces de dicha ecuación determinan la forma de la respuesta, donde el sistema es sobreamortiguado si son reales y distintas, criticamente amortiguado si son reales e iguales, y subamortiguado sin son complejas. Las raices estan dadas por la ecuación:
\[ s_{1,2} = -\alpha \pm \sqrt{\alpha^2 - \omega_0^2} \]

En esta expresión, $\alpha$ representa el factor de atenuación y $\omega_0$ la frecuencia de resonancia natural:
\[ \alpha = \frac{R_{eq}}{2L} \]
\[ \omega_0 = \frac{1}{\sqrt{LC}} \]

La relación entre estos parámetros define tres regímenes posibles. El sistema es sobreamortiguado si $\alpha > \omega_0$, críticamente amortiguado si $\alpha = \omega_0$, y subamortiguado si $\alpha < \omega_0$. En este último caso, el sistema presenta una oscilación con una frecuencia propia amortiguada $\omega_d$:
\[ \omega_d = \sqrt{\omega_0^2 - \alpha^2} \]
\newpage
